% from ega4-18-fr.tex
\oldpage[IV-4]{193}

\begin{proposition}[19.2.8]
\label{IV.19.2.8}
Soient $X$ un $S$-préschéma, $Y$ un sous-préschéma fermé de $X$ défini par un Idéal quasi-cohérent $\mathscr J$ de $\sh{O}_X$, $u:Y\to X$ l'injection canonique. Si $u$ est transversalement régulière relativement à $S$, les $\sh{O}_X$-Modules $\mathscr J^n$ et $\mathscr J^n/\mathscr J^m$ ($0\leq n<m$) sont $S$-plats aux points de $Y$. Pour tout changement de base $S'\to S$, si $X'=X_{(S')}$, $Y'=Y_{(S')}$ et $u'=u_{(S')}$, on a
\[
  \sh{G}r_n(u')=\sh{G}r_n(u)\otimes_{\sh{O}_S}\sh{O}_{S'}
\]
à isomorphisme près pour tout $n>0$.
\end{proposition}

En effet, $\sh{O}_X$ et $\sh{O}_X/\mathscr J$ sont par hypothèse $S$-plats aux points de $Y$, et $\mathscr J^n/\mathscr J^{n+1}$ est un $(\sh{O}_X/\mathscr J)$-module localement libre (16.9.3), donc est $S$-plat aux points de $Y$. La première assertion résulte alors de ($0_{\mathrm I}$, 6.1.2) et la seconde de (11.2.9.2).

La proposition suivante généralise et précise (17.16.1):

\begin{proposition}[19.2.9]
\label{IV.19.2.9}
Soient $f:X\to S$ un morphisme, $x$ un point de $X$, $s=f(x)$. Supposons que $f$ soit de présentation finie au point $x$, et soit un morphisme de Cohen-Macaulay au point $x$ (6.8.1). Alors il existe un sous-préschéma $Y$ de $X$ contenant $x$ et ayant les propriétés suivantes:
\begin{enumerate}
  \item[(i)] L'injection canonique $j:Y\to X$ est transversalement régulière relativement à $S$.
  \item[(ii)] Le morphisme $g=f\circ j:Y\to S$ est de Cohen-Macaulay (ce qui équivaut à dire, compte tenu de (i), que toutes les fibres $g^{-1}(g(y))$ de $g$ sont des préschémas de Cohen-Macaulay).
  \item[(iii)] Le point $x$ est point maximal de $Y_s=g^{-1}(s)$.
\end{enumerate}
\end{proposition}

Posons $X_s=f^{-1}(s)$; par hypothèse, l'anneau $\sh{O}_{X_s,x}$ est de Cohen-Macaulay; soit $(t_i)_{1\leq i\leq n}$ un système de paramètres de cet anneau, qui est donc une suite régulière ($0$, 16.5.7). Il y a un voisinage ouvert $U$ de $x$ dans $X$ et des sections $h_i$ ($1\leq i\leq n$) de $\sh{O}_X$ au-dessus de $U$ telles que les $t_i$ soient les images canoniques des $(h_i)_x$ dans l'anneau quotient $\sh{O}_{X_s,x}$ de $\sh{O}_{X,x}$. Soit $\mathscr J$ l'Idéal de $\sh{O}_U$ engendré par les $h_i$, et prenons pour $Y$ le sous-préschéma fermé de $U$ défini par $\mathscr J$. On peut supposer que $f|U$ est localement de présentation finie, et il en est donc de même de $g$ par définition de $Y$; il résulte alors de l'équivalence de $c)$ et $a)$ dans (11.3.8) que la suite $((h_i)_x)$ est régulière dans $\sh{O}_{X,x}$ et que $\sh{O}_{Y,x}$ est plat sur $\sh{O}_{S,s}$; en vertu de (11.3.1) $f$ et $g$ sont donc plats dans un voisinage de $x$ et il résulte de (19.2.4) que $j$ est une immersion transversalement régulière relativement à $S$ (en remplaçant au besoin $Y$ par $Y\cap V$ pour un voisinage ouvert $V$ de $x$ dans $X$). On a ainsi prouvé (i) et on voit qu'on peut supposer que $g:Y\to S$ est un morphisme plat. Comme $\sh{O}_{Y_s,x}$, quotient de $\sh{O}_{X_s,x}$ par l'idéal engendré par les $t_i$, est artinien par construction ($0$, 16.3.6), $x$ est point maximal de $Y_s$. Enfin, l'anneau $\sh{O}_{Y_s,x}$ étant artinien, est un anneau de Cohen-Macaulay ($0$, 16.5.1), donc, en remplaçant au besoin $Y$ par son intersection avec un voisinage ouvert de $x$ dans $X$, on peut supposer que $g$ est un morphisme de Cohen-Macaulay, en vertu de (12.1.7, (iii)).

En particulier, on retrouve l'existence des « quasi-sections » \emph{plates} $Y$ d'un morphisme $f:X\to S$ en un point $x\in X_s$ qui est \emph{fermé} dans $X_s$ et tel que $f$ soit plat au point $x$ et que $\sh{O}_{X_s,x}$ soit un anneau de Cohen-Macaulay (17.16.1) et on voit en outre que l'immersion $Y\to X$ est \emph{transversalement régulière} relativement à $S$. Ce dernier complément est aussi valable pour la quasi-section étale définie dans (17.16.3, (i)).

\subsection{Intersections complètes relatives (cas plat).}
\EGAReaderTarget{ega4.c04.subsec.19-3}{subsection:IV.19.3}

